\documentclass{beamer}
% definitions for slides for the semantics course - CSC501
% Lutz Hamel, (c) 2006

\usetheme{Warsaw}
\usepackage{bussproofs}
\usepackage{amsmath}
\usepackage{amssymb}
\usepackage{latexsym}
\usepackage{xypic}
\usepackage{alltt}
\usepackage{rotating}
\usepackage{graphicx}
\usepackage{url}

\newcommand{\ul}[1]{\underline{#1}}
\newcommand{\orbar}{\,|\,}
\newcommand{\co}{\,\colon\;}
\newcommand{\syntaxset}[1]{\ensuremath{\mbox{\bf #1}}}
\newcommand{\nonterm}[1]{\ensuremath{\mbox{#1}}}
\newcommand{\term}[1]{\ensuremath{\mbox{\bf #1}}}
\newcommand{\ifstmt}[3]{\ensuremath{{\bf if}\; {#1}\;{\bf then}\;{#2}\;{\bf else}\;{#3}\;\term{end}}}
\newcommand{\whilestmt}[2]{\ensuremath{{\bf while}\; {#1}\;{\bf do}\;{#2}\; \term{end}}}
\newcommand{\funcstmt}[3]{\ensuremath{{\bf fun}\; {#1}\; {\bf is}\; {#2} \; {\bf return}\; {#3}}}
\newcommand{\localstmt}[1]{\ensuremath{{\bf local}\; {#1}}}
\newcommand{\pairmap}[3]{\ensuremath{( {#1}, {#2} ) \mapsto {#3}}}
\newcommand{\single}[1]{\ensuremath{\langle {#1}\rangle}}
\newcommand{\pair}[2]{\ensuremath{({#1}, {#2} )}}
\newcommand{\triple}[3]{\ensuremath{\langle {#1}, {#2}, {#3} \rangle}}
\newcommand{\cond}[3]{\ensuremath{({#1}?\,{#2} : {#3})}}
\newcommand{\condbar}[3]{\ensuremath{({#1}\rightarrow{#2} \mid {#3})}}
\newcommand{\sem}[2]{\ensuremath{{#1}-\!\!-\!\!\!\!\gg{#2}}}
\newcommand{\liftfunc}[1]{\ensuremath{\lfloor{#1}\rfloor}}
\newcommand{\bcode}{\begin{alltt}\scriptsize}
\newcommand{\ecode}{\end{alltt}}
\newcommand{\comment}[1]{}

\newcommand{\fframe}[1]{
\begin{center}
\fbox{
\begin{minipage}{3.6in}
{#1}
\end{minipage}
}
\end{center}
}

\begin{document}

\begin{frame}{CSC 501 -- Semantics of Programming Languages}
\begin{itemize}
\item Subtitle: An Introduction to Formal Methods.
\item	Instructor: Dr. Lutz Hamel
\item	Email: lutzhamel@uri.edu
\item	Office: Tyler, Rm 251
\end{itemize}

\vspace{.1in}

\begin{center}
    \includegraphics[height=35mm]{images/hokusai}
\end{center}

\end{frame}

\begin{frame}{Books}
There are no required books in this course; however, occasionally I will assign readings based on material
available on the web.
\end{frame}

\begin{frame}{Course Objectives}

The aim of this course is to
\begin{itemize}
\item Familiarize you with the basic techniques of applying formal methods to programming languages.

\item This includes constructing models for programming languages and using these models to prove properties such as correctness and equivalence of programs.

\item Look at all major programming language constructs including assignments, loops, type systems, and procedure calls together with their models.

\item Introduce mechanical theorem provers so that we can test and prove properties of non-trivial programs.
\end{itemize}
\end{frame}

\begin{frame}{Some Definitions}
\vspace{.1in}
{\bf Definition}: In {\bf programming language semantics} we are concerned with the {\em rigorous mathematical study} of the {\em meaning} of programming languages. The meaning of a language is given by a {\em formal system} that describes the possible computations expressible within that language.
\end{frame}

\begin{frame}{Some Definitions}
\vspace{.1in}
{\bf Definition:} In computer science and software engineering, {\bf formal methods} are  techniques for the specification, development and verification of software and hardware systems based on {\em formal systems}.

\end{frame}


\begin{frame}{Formal Systems}
\vspace{.1in}
{\bf Definition:}
A {\bf formal system} consists of a {\em formal language} and a set of  {\em inference rules}.
The formal language is composed of
primitive symbols that make up well formed formulas and the inference rules are used to derive expressions from other expressions within the formal system. A formal system may be formulated and studied for its intrinsic properties, or it may be intended as a description (i.e. a model) of external phenomena.\footnote{Wikipedia}

\vspace{.1in}
In order to be truly useful in computer science, we require our formal systems to be {\em machine executable}.

\end{frame}

\begin{frame}{Uses of Formal Methods}

\begin{description}
\item[\em Implementation Issues] Formally specified models can be considered machine-independent specifications of
program behavior. They can act as ``yard sticks'' for the correctness of program implementations, transformations, and
optimizations.
\item[\em Verification] Basis of methods for reasoning about program
properties (e.g. equivalence) and program specifications (program correctness).
\item[\em Language Design] Can bring to light ambiguities and unforeseen
subtleties in programming language constructs.
\end{description}
\end{frame}





\begin{frame}[fragile]{Observations}
When programming we can observe two mental activities:
\begin{itemize}
\item We construct {\em correct looking} programs - {\em syntactically} correct programs.
\item We construct {\em models} of the intended computation in our minds. Consider,
\bcode
x := 1
while (x <= 10) do
     writeln(x)
     x := x + 1
end whiledo
\ecode
Any person with some familiarity of programming immediately has a mental picture that this program
will generate a list of integers from 1 through 10.
\end{itemize}

\end{frame}

\begin{frame}{Programming Language Definitions}

Mirroring our intuition, language definitions consist of two parts:

\begin{description}
\item[\em Syntax] The formal description of the
{\bf structure} of  well-formed expressions, phrases, programs, etc.
\item[\em Semantics] The formal description of the {\bf  meaning} of the syntactic features of a programming language
usually understood in terms of the runtime {\bf behavior} each syntactic construct evokes.  The formal description of the behavior of all
the syntactic features of a language is considered a {\bf model} of the language.
\end{description}
\end{frame}

\begin{frame}{Evaluation/Interpretation}
Syntax and semantics of a programming language are usually related via an {\em evaluation relation}
or {\em interpretation}, say $h$.  Then we say that the interpretation $h$ takes each syntactic element and
maps it into the appropriate semantic construct.

\vspace{.1in}

We often represent this with the diagram
\[
\xymatrix{
\mathit{Semantics}\\
\mathit{Syntax}\ar[u]^h
}
\]
{\bf Note:} In order for the interpretation $h$ to make any sense we will have to define the syntax and
semantics in terms of sets.

\end{frame}


\begin{frame}{Formal Systems and Programs}

The formal systems we will be using in this course are:
\begin{itemize}
\item First-order logic extended with natural deduction -- natural semantics.
\item The {\em first order predicate calculus} (often also called first order logic) to construct semantics of programming languages.
\end{itemize}
\end{frame}

\begin{frame}{Readings}
\begin{itemize}
\item Read Chapter 0 in "Denotational Semantics" by David Schmidt (available from the course website).
\item Read Sections 2.1 and 2.2 in "Denotational Semantics" by David Schmidt.
\end{itemize}
\end{frame}

\begin{frame}{\large Mathematical  Preliminaries: Logical Notation}

We\footnote{\tiny The material presented here is based on ``Naive Set Theory'' by P. Halmos and ``The Formal Semantics of Programming Languages'' by G. Winskel.}
will use first order logic as the basis for our reasoning.  Without going into the formal details
of first order logic terminology and sentence construction we have the following statements:
\begin{itemize}
\item $A \wedge B$ denotes the conjunction $A$ and $B$,
\item $A \vee B$ denotes the disjunction $A$ or $B$,
\item $\neg A$ denotes the negation not $A$,
\item $A \Rightarrow B$ denotes the implication, if $A$ then $B$,
\item $A \Leftrightarrow B$ denotes the logical equivalence, $A$ if and only if $B$ (often
written as $A$ iff $B$),
\end{itemize}
where $A$ and $B$ are statements or assertions.
\vspace{.1in}
 \end{frame}

\begin{frame}{\large Mathematical  Preliminaries: Logical Notation}
Observe the precedences of the logical operators, ordered from high to low:
\begin{itemize}
\item $\neg$
\item $\wedge,\vee$
\item $\Rightarrow, \Leftrightarrow$
\end{itemize}
\end{frame}


\begin{frame}{A Word or Two about Implication}
The truth table for the implication operator `$\Rightarrow$' can be given as
{\scriptsize
\[
\begin{array}{lcc|c}
&A & B & A \Rightarrow B\\ \hline
(1)& 1 & 0 & 0\\
(2) & 1 & 1 & 1\\
(3) & 0 & 0 & 1\\
(4) & 0 & 1 & 1
\end{array}
\]
}
Entries $(1)$ and $(2)$ are intuitive: When the antecedent $A$ is true but the consequent $B$ is false then
the implication itself is false.  If both the antecedent and the consequent are true then the implication is true.

\vspace{.1in}

However, entries $(3)$ and $(4)$ are somewhat counter intuitive.  They state that if the antecedent $A$ is false
then the implication is true regardless of the value of the consequent.  In other words, we can conclude ``anything''
from an antecedent that is false.  In mathematical jargon we say that $(3)$ and $(4)$  {\bf hold trivially}.
\end{frame}

\begin{frame}{A Word or Two about Implication -- An Example}

{\scriptsize
\[
\begin{array}{l}
\mbox{If Bob is a bachelor, then he is single.}\\
\mbox{Bob is a bachelor.}\\ \hline
\therefore\mbox{Bob is single.}
\end{array}
\]
}
Now consider an antecedent that is not true,
{\scriptsize
\[
\begin{array}{l}
\mbox{If Bob is a bachelor, then he is single.}\\
\mbox{Bob is not a bachelor.}\\ \hline
\therefore\mbox{Bob is not single (by rule $(3)$).}
\end{array}
\]
}
Since the antecedent is not true rule $(3)$ allows us to conclude the
opposite of what the implication dictates.  However, the following is also
valid reasoning,
{\scriptsize
\[
\begin{array}{l}
\mbox{If Bob is a bachelor, then he is single.}\\
\mbox{Bob is not a bachelor.}\\ \hline
\therefore\mbox{Bob is  single (by rule $(4)$).}
\end{array}
\]
}
Not being a bachelor does not necessarily imply that Bob is not single.  For example,
Bob could be a widower or a divorcee.
\end{frame}

\begin{frame}{A Word or Two about Implication}

Given the truth table for implication,
\[
\begin{array}{lcc|c}
&A & B & A \Rightarrow B\\ \hline
(1)& 1 & 0 & 0\\
(2) & 1 & 1 & 1\\
(3) & 0 & 0 & 1\\
(4) & 0 & 1 & 1
\end{array}
\]
this means that in order to show that an implication holds we only have to show that rule $(2)$
holds.  Rule $(1)$ states that the implication is false and rules $(3)$ and $(4)$ are trivially true and
therefore not interesting.
\end{frame}

\begin{frame}{Closely Related: Equivalence}

We write $A \Leftrightarrow B$ if $A$ and $B$ are equivalent.

Given the truth table for the equivalence operator is given as,
\[
\begin{array}{lcc|c}
&A & B & A \Leftrightarrow B\\ \hline
(1)& 1 & 0 & 0\\
(2) & 1 & 1 & 1\\
(3) & 0 & 0 & 1\\
(4) & 0 & 1 & 0
\end{array}
\]
That is, the operator only produces a true value if $A$ and $B$ have the same truth assignment.

Another, and very useful, way to look at the equivalence operator is as follows:
\[
A \Leftrightarrow B \equiv A \Rightarrow B \wedge B \Rightarrow A
\]
{\bf Exercise:} Construct  the above truth table using this definition of the equivalence operator.
\end{frame}

\begin{frame}{\large Mathematical  Preliminaries: Logical Notation}
We also allow predicates (properties) as part of our notation,
\[
P(x)
\]
where the predicate $P$ is true if it holds for $x$ otherwise it is false.  We view our
standard relational operators as binary predicates.  For example, the predicate
$P(x)$ that expresses the fact that $x$ is less or equal to 3 is written as,
\[
P(x) \equiv x \le 3.
\]
{\bf Note:} Predicates can have arities larger than 1, e.g. $P(x,y)$ with $P(x,y)\equiv x \le y$.
\end{frame}

\begin{frame}{\large Mathematical  Preliminaries: Logical Notation}
We also allow for the quantifiers $\exists$ (there exists) and $\forall$ (for all) in our logical
statements,
\begin{itemize}
\item $\exists x.\,P(x)$ -- ``there exists an $x$ such that $P(x)$''
\item $\forall x.\,P(x)$ -- ''for all $x$ such that $P(x)$''
\end{itemize}
Some examples,
\begin{itemize}
\item $\forall x,\exists y.\, y = x^2$
\item $\forall x, \forall y.\, \mbox{female}(x) \wedge \mbox{child}(x,y) \Rightarrow \mbox{mother}(x,y)$
\end{itemize}

\end{frame}


\begin{frame}{\large Mathematical  Preliminaries: Sets}
Sets\footnote{\tiny Read Sections 2.1 and 2.2 in the book by David Schmidt.}
 are unordered collections of objects and are usually denoted by capital letters.  For example,
let $a,b,c$ denote some objects then the set $A$ of these objects is written as,
\[
A = \{a,b,c\}.
\]
There are a number of standard sets which come in handy,
\begin{itemize}
\item $\emptyset$ denotes the empty set, i. e. $\emptyset = \{ \}$,
\item $\mathbb N$ denotes the set of all natural numbers including 0, e. g. ${\mathbb N} = \{0,1,2,3,\cdots\}$,
\item $\mathbb I$ denotes the set of all integers, ${\mathbb I}= \{\cdots,-2,-1,0,1,2,\cdots\}$,
\item $\mathbb R$ denotes the set of all reals,
\item $\mathbb B$ denotes the set of boolean values, ${\mathbb B} = \{\mathit{true}, \mathit{false} \}$.
\end{itemize}
\end{frame}

\begin{frame}{\large Mathematical  Preliminaries: Sets}
The most fundamental property in set theory is the notion of {\em belonging},
\[
a \in A \mbox{ iff  $a$ is an element of the set $A$}.
\]
The notion of belonging allows us
to define {\em subsets},
\[
Z \subseteq A \mbox{ iff } \forall e\in Z.\, e\in A.
\]
We define set {\em equivalence} as,
\[
A = B \mbox{ iff } A \subseteq B \wedge B \subseteq A
\]
\end{frame}


\begin{frame}{\large Mathematical  Preliminaries: Sets}
We can construct new sets from given sets using {\em union},
\[
A \cup B = \{ e \mid e\in A \vee e \in B\},
\]
 and {\em intersection},
\[
A \cap B = \{ e \mid  e\in A \wedge e \in B\}.
\]
There is another important set construction called the {\em cross product},
\[
A\times B = \{ (a,b) \mid  a\in A \wedge b\in B\},
\]
$A\times B$ is the set of all ordered pairs where the first component of the pair is drawn from
the set $A$ and the second component of the pair is drawn from $B$. (

{\bf Exercise:} Let $A=\{a,b\}$ and $B=\{c,d\}$, construct the set $A\times B$.

\end{frame}

\begin{frame}{\large Mathematical  Preliminaries: Sets}
A construction using subsets is the {\em powerset} of some set $X$,
\[
{\mathcal P}(X),
\]
The {\em powerset} of set $X$ is set of all subsets of $X$.  For example,
let $X = \{a,b\}$, then
\[
{\mathcal P}(X) = \{ \emptyset, \{a\}, \{b\}, \{a,b\} \}.
\]
{\bf Note:} $\emptyset \subset X$

\vspace{.1in}

{\bf Exercise:} What would ${\mathcal P}(X\times X)$ look like?

\end{frame}

\begin{frame}{\large Mathematical  Preliminaries: Sets}
\scriptsize
The fact that $\emptyset \subset X$ for any set $X$ is interesting in its own right.  Let's see if we can prove it.

\vspace{.1in}

{\bf Proof:} Proof by contradiction. Assume $X$ is any set. Assume that $\emptyset$ is not a subset of $X$.
Then the definition of subsets,
\[
A \subseteq B \Leftrightarrow \forall e \in A. e \in B,
\]
implies that there exist at least one element in $\emptyset$ that is not also in $X$.  But that is not possible
because $\emptyset$ has no elements -- a contradiction.  Therefore, our assumption the
that $\emptyset$ is not a subset of $X$ must be wrong and we can conclude that $\emptyset \subset X$.

\end{frame}


\begin{frame}{\large Mathematical  Preliminaries: Relations}
A (binary) relation is a set of ordered pairs.  If $R$ is a relation that relates the elements of
set $A$ to the elements $B$, then
\[
R \subseteq A\times B.
\]
This means if $a\in A$ is related to $b\in B$ via the relation $R$, then $(a,b)\in R$.  We often
write
\[
a\, R\, b.
\]
Consider the relational operator $\le$ applied to the set ${\mathbb N}\times {\mathbb N}$.
This induces a relation, call it $\le \subseteq {\mathbb N}\times {\mathbb N}$, with
$(a,b) \in \le$ (or $a \le b$ in our relational notation) if $a\in{\mathbb N}$ is less or equal to $b\in{\mathbb N}$.
\end{frame}

\begin{frame}{\large Mathematical  Preliminaries: Relations}
The first and second components of each pair in some relation $R$ are drawn from different sets called
the {\em projections} of $R$ onto the first and second {\em coordinate}, respectively.  We introduce
the operators {\em domain} and {\em range} to accomplish these projections.  Let $R\subseteq A\times B$,
then,
\[
\mbox{dom}(R) = A,
\]
and
\[
\mbox{ran}(R) = B.
\]
In this case we talk about a relation {\em from} $A$ {\em to} $B$.  The range is often
called the co-domain. If $R\subseteq X\times X$, then
\[
\mbox{dom}(R) = \mbox{ran}(R) = X.
\]
Here we talk about a relation {\em in} $X$.
\end{frame}

\begin{frame}{\large Mathematical  Preliminaries: Relations}
Let $R \subseteq X\times X$ such that $(a,b)\in R$ iff $a=b$.  That is, $R$ is the {\em equality relation}
in  $X$. (What do the elements of the equality relation look like for $\mathbb{N}\times\mathbb{N}$?)

\vspace{.1in}

A relation $R\subseteq X\times X$ is an {\em equivalence relation} if the following conditions hold,
\begin{itemize}
\item $R$ is {\em reflexive}\footnote{Recall that $x\, R\, x \equiv (x,x)\in R$} -- $x\, R\, x$,
\item $R$ is {\em symmetric} -- $x\, R\, y \Rightarrow y\, R\, x$,
\item $R$ is {\em transitive} -- $x\, R\, y \wedge y\, R\, z \Rightarrow x\, R\, z$,
\end{itemize}
where $x,y,z \in X$.

\vspace{.1in}

The {\em smallest} equivalence relation in some set $X$ is the equality relation defined above.
The {\em largest} equivalence relation is some set is the cross product $X\times X$. (Consider the
smallest/largest equiv. relation in $\mathbb I$)
\end{frame}

\begin{frame}{\large Mathematical  Preliminaries: Functions}
A {\em function} $f$ from $X$ to $Y$ is a relation $f \subseteq X\times Y$ such that
\[
\forall x\in X, \exists y,z\in Y.\, (x,y)\in f \wedge (x,z)\in f \Rightarrow y = z.
\]
In other words, each $x\in X$ has a unique value $y\in Y$ with $(x,y)\in f$ or functions are constrained relations.

\vspace{.1in}

We let $X \rightarrow Y$ denote the {\em set of all functions} from $X$ to $Y$
(i. e. $X \rightarrow Y \subset {\mathcal P}(X\times Y)$, why is the subset strict? Hint: it is not a relation), then the customary notation
for specifying functions can be defined as follows,
\[
f\co X \rightarrow Y \mbox{ iff } f\in X\rightarrow Y.
\]
\end{frame}

\begin{frame}{\large Mathematical  Preliminaries: Functions}
For {\em function application} it is customary to write
\[
f(x) = y
\]
for $(x,y)\in f$.  In this case we say that the function is {\em defined} at point $x$.
Otherwise we say that the function is {\em undefined} at point $x$ and we write $f(x)=\perp$.

\vspace{.1in}

Note that $f(\perp) = \perp$  and we say the $f$ is {\em strict}.
\vspace{.1in}

We say that $f\co X \rightarrow Y$ is a {\em total} function if $f$ is defined for all $x\in X$.
Otherwise we say that $f$ is a {\em partial} function.
\end{frame}

\begin{frame}{\large Mathematical  Preliminaries: Functions}
We can now make the notion of a predicate formal -- a predicate is a function whose range (co-domain) is restricted to the
boolean values:
\[
P: X \rightarrow {\mathbb B}
\]
where $P$ is a predicate that returns true or false for the objects in set $X$.

\vspace{.1in}

{\bf Example:} Let $U$ be the set of all possible objects -- a universe if you like, and let,
\[
\mathit{human}: U \rightarrow {\mathbb B}
\]
be the predicate that returns true if the object is a human and will return false otherwise, then
\begin{eqnarray*}
\mathit{human}(\mbox{socrates}) = \mathit{true}\\
\mathit{human}(\mbox{car}) = \mathit{false}
\end{eqnarray*}
\end{frame}


\begin{frame}{\large Mathematical  Preliminaries: Exercises}
\begin{enumerate}
\item In your own words explain what the function $m\co X\times Y \rightarrow Z$ does.
\item How would you describe the function $c \co X \rightarrow (Y \rightarrow Z)$?
\item In your own words explain what the relation $R \subseteq (X\times Y)\times (Z\times W)$ does.
\end{enumerate}
\end{frame}





\end{document}
%%%%%%%%%%%%%%%%%%%%%%%%%%% end of template1.tex %%%%%%%%%%%%%%%%%%%%%%%%%%%%%%%%
